\documentclass[a0,landscape]{a0poster}

\usepackage[paper=a0paper, landscape, margin=2cm]{geometry}

\usepackage{multicol} % This is so we can have multiple columns of text side-by-side
\columnsep=100pt % This is the amount of white space between the columns in the poster

\usepackage{tikz}

\usepackage{sectsty}
\sectionfont{\huge\selectfont}

\usepackage[scaled]{helvet}
\renewcommand\familydefault{\sfdefault} 
\usepackage[T1]{fontenc}

\usepackage{graphicx} % Required for including images
\graphicspath{{figures/}} % Location of the graphics files
\usepackage{booktabs} % Top and bottom rules for table
\usepackage[font=small,labelfont=bf]{caption} % Required for specifying captions to tables and figures
\usepackage{amsfonts, amsmath, amsthm, amssymb} % For math fonts, symbols and environments
\usepackage{wrapfig} % Allows wrapping text around tables and figures

\begin{document}

\tikz[remember picture,overlay] \node[inner sep=0pt, anchor=north] at (current page.north){\includegraphics[width=\paperwidth,height=12cm]{GreenBackground.pdf}};

\vspace{-1cm}
\begin{minipage}[t]{0.3\linewidth}
\vspace{0pt}
\includegraphics[width=0.9\textwidth]{UofALogo.pdf}
\end{minipage}
%
\begin{minipage}[t]{0.7\linewidth}
\vspace{0.75cm}
\color{white}
\veryHuge \textbf{\textsc{Redliner: An Activity Monitor for Manual Wheelchair Users}}\\[1cm]
\huge {Kenton Hamaluik, MSc, EIT\textsuperscript{1}, Zohreh Salimi, MSc\textsuperscript{1}, Martin Ferguson-Pell, PhD, C.Phys.\textsuperscript{1,2}}\\
\large \textsuperscript{1}Faculty of Rehabilitation Medicine,\ \textsuperscript{2}Peter Lougheed Leadership College\\
\end{minipage}

\vspace{4cm}

\begin{multicols*}{3}
\large

%\begin{abstract}
%Activity monitors such as those produced by FitBit and Garmin are popular amongst the general population, however despite their usefulness in healthy living they are unsuitable for wheelchair users as they do not measure wheelchair propulsion. While the SmartWheel is an existing wheelchair activity monitor, it is prohibitively expensive and intrusive for personal use. The objective of this work was to create a low-cost manual wheelchair activity monitor capable of measuring parameters relating to activity level while also monitoring and notifying the user of damaging over-exertion events, fulfilling the role of personal activity monitor for wheelchair users. The Redliner device uses a novel non-intrusive sensor which attaches to a user's wheel and measures the wheelchair's motion. A wheelchair-specific propulsion algorithm was developed which counts pushes, estimates distance travelled and time spent active, and enumerates overexertion events in real-time onboard the device. In order to evaluate the effectiveness of the Redliner device, a series of trials collecting a minimum of 3 pushes on various surfaces including linoleum, grass, gravel, and rough pavement were performed while collecting synchronous SmartWheel data. These surfaces were chosen to represent real-world wheelchair propulsion situations which could prove difficult for the device. Kinematic data produced by the Redliner device was found to be in good agreement with the SmartWheel, with only slight losses in accuracy on extremely rough surfaces. Parameters reported by the propulsion algorithm closely matched experimental conditions. Redliner is a novel device for wheelchair users which has been proven effective for activity and over-exertion monitoring.
%\end{abstract}

\section*{Background}
\begin{itemize}
    \item Wheelchair users have a strong need to monitor activity
        \begin{itemize}
            \item Manual wheelchair users are highly susceptible to repetitive stress injuries \cite{big}
            \item Propulsion forces over 80\% of maximum capacity result in damage \cite{small}
        \end{itemize}
    \item Typical existing activity monitors don't work with manual wheelchair users
    \item SmartWheels are the gold standard, but are an expensive clinical tool
        \begin{itemize}
            \item There is a need for an affordable consumer-grade tool for activity monitoring
        \end{itemize}
\end{itemize}

\section*{Objective}
The objective of this work was to create an inexpensive activity monitor for manual wheelchair users capable of measuring the following data:

\begin{itemize}
    \item Number of propulsion strokes
    \item Average travel velocity
    \item Amount of time spent active
    \item Estimated distance travelled
    \item Number of ``redline events''\textsuperscript{1}
\end{itemize}

\textsuperscript{1} \textit{redline events} are instances of when the user's propulsion force exceeds 80\% of the maximal propulsion force they can generate

\section*{Method}

\section*{Result}

\section*{Conclusions}
\begin{itemize}
    \item Redliner is a new activity monitor for manual wheelchair users
    \item Redliner has been validated against expensive SmartWheel devices
    \item Redliner is moving forward as a commercial entity to produce and sell the devices
\end{itemize}

\nocite{*} % Print all references regardless of whether they were cited in the poster or not
\bibliographystyle{plain} % Plain referencing style
\bibliography{sources} % Use the example bibliography file sample.bib

\section*{Acknowledgements}
Funding for this work was provided for by the University of Alberta, Telus, the Government of Canada, and Redliner Inc.

\end{multicols*}
\end{document}
